\section{Related Work}
\label{sec:related}

\subsection{Numerical Weather Prediction and Regional Forecasting}
Operational weather forecasting has historically relied upon solving the discretized Navier-Stokes equations along with thermodynamics and continuity equations via high-performance computing clusters~\cite{bauer2015quiet, lynch2008origins}. The ECMWF Integrated Forecasting System (IFS HRES) represents the benchmark global numerical model. To obtain higher resolution over localized regions, regional NWP models such as the Weather Research and Forecasting (WRF) model~\cite{skamarock2008description} are typically nested inside global boundary conditions. However, regional NWP demands extensive expertise in parameterizing boundary-layer processes, microphysics, and radiation schemes, while remaining computationally demanding and prone to lateral boundary error propagation.

\subsection{Deep Learning for Global Weather Forecasting}
The integration of deep learning into weather forecasting initially started with convolutional neural networks (CNNs) and recurrent architectures~\cite{weyn2020improving, rasp2020weatherbench}. A monumental leap occurred between 2022 and 2023 with FourCastNet~\cite{pathak2022fourcastnet}, which utilized Adaptive Fourier Neural Operators (AFNO) to capture global dependencies in the spectral domain. Concurrently, Pangu-Weather~\cite{bi2023accurate} leveraged a 3D Earth-Specific Transformer architecture with hierarchical downsampling and temporal patch aggregations, demonstrating that machine learning can surpass operational IFS across major deterministic scorecard metrics. GraphCast~\cite{lam2023learning} formulated global forecasting as an Encoder--Processor--Decoder message-passing network over an icosahedral spherical mesh, setting benchmark scores in the ECMWF Scorecard. Other notable first-wave architectures include FuXi~\cite{chen2023fuxi}, FengWu~\cite{chen2023fengwu}, and ClimaX~\cite{nguyen2023climax}.

\subsection{Atmospheric Foundation Models and Generative Ensembles}
The 2024--2025 literature witnessed a paradigm shift toward massive multi-source foundation models and generative diffusion architectures. Microsoft introduced Aurora~\cite{bodnar2024aurora}, a 1.3-billion parameter atmospheric foundation model pre-trained on diverse geophysical analyses and satellite records, scaling zero-shot fine-tuning across extreme storms at $0.1^\circ$ resolution with 90\% compute savings. ECMWF developed and operationalized AIFS~\cite{lang2024aifs}, combining spherical GNN encoders with sliding-window Transformer processors directly into operational numerical forecasting. To resolve the spectral blurring and conditional mean collapse intrinsic to deterministic L2/MSE models, Google DeepMind introduced GenCast~\cite{price2024gencast}, a conditional score-based diffusion model on spherical icosahedral meshes that generates realistic 50-member ensemble trajectories outperforming operational ensemble systems (ECMWF ENS) up to 15 days. Concurrently, NASA and IBM launched Prithvi WxC~\cite{schmude2024prithvi}, a 2.3-billion parameter model for regional downscaling, while FengWu-GHR~\cite{han2024fengwu_ghr} pushed global deterministic resolution to $0.09^\circ$ ($\approx 9$~km), and FuXi-Extreme~\cite{zhong2024fuxi_extreme} leveraged diffusion to capture localized convective rainfall extremes. NeuralGCM~\cite{kochkov2024neuralgcm} established a hybrid modeling frontier by coupling a differentiable dynamical core (solving Navier-Stokes equations) with learned neural parameterizations for sub-grid processes.

\subsection{The 2025--2026 Frontier: Diffusion LAMs, Latent Consistency, and Tropical GNNs}
Entering 2025 and 2026, research has sharply pivoted from compute-heavy global foundation models toward ultra-fast generative architectures, regional limited area models (LAMs), and kilometric downscaling:
\begin{itemize}
    \item \textbf{Diffusion-LAM (2025):} Larsson et~al.~\cite{larsson2025diffusion_lam} formulated the first dedicated regional limited-area model driven by conditional diffusion. By buffering regional predictions with time-varying lateral boundary conditions from surrounding reanalysis, Diffusion-LAM eliminates unphysical boundary shocks while retaining fine-scale kinetic energy spectra without high-frequency smoothing.
    \item \textbf{PuYun-LDM (2026):} Wu et~al.~\cite{wu2026puyun_ldm} introduced a high-resolution Latent Diffusion Model combining a 3D Masked AutoEncoder (3D-MAE) with Variable-Aware Masked Frequency Modeling (VA-MFM). By operating in a compact latent subspace, PuYun-LDM generates 15-day global ensemble forecasts at $0.25^\circ$ in only 5 minutes on a single GPU, dramatically slashing generative computational barriers.
    \item \textbf{AirCast-SR (2026):} Addressing the persistent resolution gap in tropical complex terrain, Luitel et~al.~\cite{luitel2026aircast_sr} deployed Latent Consistency Models (LCM) to achieve one-step kilometric super-resolution downscaling from $0.25^\circ$ ($\approx 28$~km) to 1~km grids, preserving meso-$\gamma$ convective structures and exhibiting zero-shot cross-basin generalization.
    \item \textbf{Spatio-Temporal GNNs for Cyclones (2025):} Architectures such as StormNet~\cite{chen2025stormnet} integrate multi-relational graph convolution with spatial cross-attention to predict tropical cyclone track deflection, central pressure deepening, and rapid intensification (RI) across the South China Sea and Western Pacific basins.
\end{itemize}

\subsection{Graph Neural Networks in Atmospheric Dynamics}
Atmospheric flow exhibits non-Euclidean spherical geometry, irregular geographical features, and multiscale teleconnections that pose challenges for standard 2D Euclidean CNNs~\cite{bronstein2017geometric}. Standard convolutions on latitude-longitude grids suffer from severe polar distortion (coordinate convergence near the poles) and assume translation invariance that conflicts with distinct geographical boundary conditions (e.g., land vs. ocean). GNNs overcome these limitations by modeling atmospheric nodes as arbitrary graph topologies where edges represent physical spatial proximity or long-range dynamical interactions~\cite{battaglia2018relational, keisler2022forecasting}. By separating the observation grid from the computational latent mesh, GNNs can achieve linear computational complexity $O(N)$ with respect to node count, making them attractive for high-resolution modeling.

\subsection{Regional AI Weather Forecasting and Multimodal Fusion}
While global models predict the entire planet, regional data-driven forecasting focuses computational budgets on local geographic sectors of interest. Recent efforts have examined regional Fourier operators~\cite{bonev2023spherical} and downscaling networks~\cite{vandal2017deepsd}. In Southeast Asia, local convective dynamics are notoriously difficult to predict from reanalysis alone because reanalysis models parameterize sub-grid precipitation. Satellite remote sensing—particularly geostationary infrared brightness temperature ($T_b$) from the NASA Global Precipitation Measurement (GPM) MERGIR dataset~\cite{janowiak2001mergir}—provides observational proxies for convective cloud-top height and convective vigor at 4 km resolution. Integrating observational satellite data alongside numerical reanalysis offers a dual-stream perspective that bridges the gap between massive global foundation models and accessible regional operational deployment on commodity hardware.
